\subsection*{E(iv). The exact quartic local-order law}

\begin{statement}{E.13}{Proved}{The full CM image and the affine image of every primitive class}
Let $E_+ /\mathbb Q$ be $y^2=x^3+x$, let $\ell\in\{3,5\}$,
put $V=T_\ell E_+$ and $K=\mathbb Q(i)$, and write
$\rho:G_{\mathbb Q}\to\operatorname{GL}(V)$ for the actual Tate action.
Then
\[
 \rho(G_K)=C_\ell,\qquad
 \rho(G_{\mathbb Q})=N_\ell=C_\ell\rtimes\langle c\rangle,
 \tag{E13.1}
\]
where $C_5=(\mathbb Z_5^\times)^2$ acts on the two integral CM
lines, $C_3=\mathbb Z_3[i]^\times$ acts on the rank-one
$\mathbb Z_3[i]$-module $V$, and $c$ interchanges the two lines at
five and acts by conjugation at three.  This is the full normalizer
of the displayed Cartan subgroup in $\operatorname{GL}(V)$.

Every continuous $V$-valued cocycle factors through
$G_{\mathbb Q,S}$ for $S=\{2,\ell,\infty\}$.
The module $H^1(G_{\mathbb Q},V)$ is finite free over
$\mathbb Z_\ell$ of positive rank.  For every primitive class in
this module and every continuous cocycle $h$ representing it, the
actual homomorphism
\[
 G_{\mathbb Q}\longrightarrow V\rtimes N_\ell,
 \qquad g\longmapsto(h(g),\rho(g))
 \tag{E13.2}
\]
is surjective.  Here primitive means not divisible by $\ell$ in
$H^1(G_{\mathbb Q},V)$.  No class-number calculation or assertion
that this cohomology module has rank one is needed for these conclusions.
\end{statement}
\begin{proof}
The endomorphism $I:(x,y)\mapsto(-x,iy)$ is defined over $K$
and satisfies $I^2=-1$.  Its determinant on a Tate module is one,
by the Weil pairing.  At five its two eigenvalues are therefore
the distinct roots of $T^2+1$, each with multiplicity one; the
projectors $(1\pm I/i_5)/2$ are integral.  At three the same
polynomial is irreducible modulo three, and $V$ is free of rank
one over $\mathbb Z_3[i]$.  Thus the centralizer of $I$ is $C_\ell$.
For $g\in G_{\mathbb Q}$ the actual endomorphism relation is
\[
       \rho(g)I\rho(g)^{-1}=I\quad(g\in G_K),\qquad
       \rho(g)I\rho(g)^{-1}=-I\quad(g\notin G_K).
 \tag{E13.3}
\]
This places the image in the claimed normalizer and identifies
its quotient of order two with $\operatorname{Gal}(K/\mathbb Q)$.
To check that the normalizer named here is the full matrix
normalizer, the units of $\mathbb Z_\ell[i]$ generate that
quadratic etale algebra: at three use $1,i$, and at five use
$(1,1)$ and $(2,1)$ in its split realization.
A matrix normalizing the unit group therefore normalizes this
algebra.  Its induced algebra automorphism is the identity or
conjugation; the kernel is the unit centralizer $C_\ell$.
Both automorphisms are represented, giving exactly $N_\ell$.

We prove the full integral Cartan image by local inertia.
At five take the Gaussian primes $\pi=1+2i$ and
$\bar\pi=1-2i$.  At the place $\pi$ of $K$, the completion is
$\mathbb Q_5$ and $\bar\pi$ is a unit.  The curve has good
reduction, since its discriminant is $-64$.  The isogeny
$[\bar\pi]$ extends to the smooth proper elliptic model and its
differential is multiplication by the unit $\bar\pi$.
Its kernel and the kernels of all its iterates are consequently
finite etale over the valuation ring.  Inertia acts trivially on
their Tate line.  Bezout for the coprime Gaussian ideals
$(\pi)$ and $(\bar\pi)$ gives the actual decomposition
\[
 E_+[5^n]=E_+[\pi^n]\oplus E_+[\bar\pi^n].
 \tag{E13.4}
\]
The determinant of inertia on these two factors is cyclotomic.
The $5$-power cyclotomic extensions of $\mathbb Q_5$ are totally
ramified, with full unit Galois group.  Hence this inertia image
is $\mathbb Z_5^\times\times\{1\}$ in one ordering of (E13.4).
Complex conjugation supplies $\{1\}\times\mathbb Z_5^\times$.
Both inertia groups lie in $G_K$, proving $\rho(G_K)=C_5$.

At three the completion $K_3=\mathbb Q_3(i)$ is the unramified
quadratic extension, with residue field of order nine.  The
reduction of $E_+$ satisfies the literal endomorphism identity
\[
          [-3](x,y)=(x^9,y^9).
 \tag{E13.5}
\]
For an elementary verification, in characteristic three the
division polynomials give $\psi_3=-1$, $\psi_2=-y$ and
$\psi_4=y(x^6+2x^4+x^2+2)$.  Therefore
\[
 x([3]P)=x+(x^3+x)(x^6+2x^4+x^2+2)=x^9.
\]
The curve equation makes the other coordinate either $y^9$
or $-y^9$ as a rational function.  The point $(2,1)$ has double
$(0,0)$ and order four over $\mathbb F_3$, so $[3](2,1)=(2,-1)$;
the sign is negative.  This proves (E13.5), without a choice of
a Frobenius eigenvalue.

In the formal parameter $T=-x/y$, the actual series
$f(T)=[-3](T)$ has derivative $-3$ and reduces to $T^9$.
It is a Lubin--Tate series for the uniformizer $-3$ of $K_3$.
The uniqueness statement for a formal group commuting with
$f$ identifies the actual formal group of $E_+$ with the
Lubin--Tate formal group, using the identity parameter map.
The same uniqueness for an endomorphism with prescribed linear
coefficient identifies each actual $[a+bi]$ with the
Lubin--Tate endomorphism for $a+bi$; continuity gives this
identification over all of $\mathbb Z_3[i]$.
Equation (E13.5) also shows that every geometric $3$-power
torsion point reduces to the origin, so all such torsion is in
this formal group.  Lubin--Tate's theorem now gives, at each
level, the totally ramified extension with Galois action
\[
      (\mathbb Z_3[i]/3^n)^\times
             \quad\hbox{on}\quad E_+[3^n].
\]
Its inverse limit is the full $C_3$, already obtained from local
inertia.  The precise formal-group input is
\cite[Lemma 1, Theorem 1, and Theorem 2(b),(d),(e)]{E_density_LT}.
In particular (E13.1) follows at both primes.  Choosing one basis
vector on a CM line and its conjugate gives the stated $c$ at
five.  At three choose a primitive vector in the $+1$ eigenspace
of $c$; it is an integral $\mathbb Z_3[i]$-basis and gives the
stated conjugation action.  These changes of basis are actual
linear isomorphisms of the Tate lattices.

To prove the assertions about cocycles, let $p\notin S$.
The Tate action on inertia $I_p$ is trivial.  A cocycle restricted
to $I_p$ is a homomorphism to the pro-$\ell$ group $V$, hence
factors through the $\ell$-primary tame quotient.
For its value $v$ on a tame generator, Frobenius equivariance
gives $(F_{E,p}-p)v=0$.  But
\[
    \det(F_{E,p}-p)=p\,\#E_+(\mathbb F_p)\ne0.
\]
Thus $v=0$.  The cocycle vanishes on all the corresponding
conjugate inertia groups and on their closed normal closure;
the action also vanishes there.  It factors through $G_{\mathbb Q,S}$.

The finite-coefficient finiteness theorem and the global Euler
characteristic formula give the following facts, with their
hypotheses satisfied since $S$ is finite and contains $\ell$
and the archimedean place:
\cite[I.4.10(c), I.4.15, and I.5.1]{E_density_MilneADT}.
The groups $H^j(G_{\mathbb Q,S},V/\ell^n)$ are finite;
they vanish for $j>2$ because $\ell$ is odd.
Their finite inverse systems have zero derived inverse limit.
The continuous groups $H^1(G_{\mathbb Q,S},V)$ and
$H^2(G_{\mathbb Q,S},V)$ are therefore compact; their reductions
modulo $\ell$ are finite, so compact Nakayama makes them finitely
generated over $\mathbb Z_\ell$.
The residual representation has no invariant vector: at three
multiplication by $i$ has irreducible polynomial, and at five
the full diagonal subgroup and its interchange have no common
fixed vector.  The coefficient exact sequence then proves that
$H^1(G_{\mathbb Q,S},V)$ has no $\ell$-torsion.
Write $r$ and $s$ for the ranks of $H^1$ and $H^2$ respectively.
Complex conjugation has one eigenvalue of each sign, and the
finite Euler formula is
\[
 \frac{\#H^1(G_{\mathbb Q,S},V/\ell^n)}
      {\#H^2(G_{\mathbb Q,S},V/\ell^n)}=\ell^n.
\]
The exact coefficient sequences give
\[
\begin{split}
0&\longrightarrow H^1(V)/\ell^n
 \longrightarrow H^1(V/\ell^n)
 \longrightarrow H^2(V)[\ell^n]\longrightarrow0,\\
H^2(V)/\ell^n&\ \xrightarrow{\sim}\ H^2(V/\ell^n),
\end{split}
\]
where these groups have source $G_{\mathbb Q,S}$.
The torsion contributions of $H^2(V)$ cancel in the ratio.
Consequently $r-s=1$, and $r\ge1$.  Primitive classes exist.

Fix such an $h$.  Its reduction $[h_1]$ in
$H^1(G_{\mathbb Q,S},V/\ell)$ is nonzero by the same exact
coefficient sequence.  Let $\Gamma_n=\rho_n(G_{\mathbb Q,S})$
and $J_n=\ker\rho_n$.  The scalar $-I$ lies in $\Gamma_n$.
For a cocycle $a:\Gamma_n\to V/\ell$, commuting $-I$ with $g$
in the cocycle identity gives
\[
       -2a(g)=(g-1)a(-I).
\]
Since two is a unit, every such $a$ is a coboundary.
Inflation--restriction shows that $h_1|_{J_n}$ is nonzero.
The subgroup $h_n(J_n)\subset V/\ell^n$ is stable under
$\Gamma_n$, and its reduction is nonzero.  This reduction is
the whole $V/\ell$: at three the CM action is irreducible;
at five a stable line for the diagonal subgroup is a coordinate
line and complex conjugation interchanges the two lines.
Nakayama therefore gives $h_n(J_n)=V/\ell^n$.
The image of $(h_n,\rho_n)$ has the full translation kernel
and surjects onto $\Gamma_n$, so is
$(V/\ell^n)\rtimes\Gamma_n$.
The image of (E13.2) is compact and surjects onto every one
of these finite quotients.  It is the entire inverse limit.
\end{proof}

\begin{statement}{E.14}{Proved}{Exact natural densities of all local Kummer orders}
Retain $E_+$, $\ell$ and any primitive cocycle $h$ of E.13.
For every rational prime $p\ne2,\ell$, choose an arithmetic
Frobenius and define
\[
 u_p=(F_{E,p}-1)^{-1}h(F_p),\qquad
 \operatorname{ord}(u_p+V)=\ell^{m_p},\quad m_p\ge0.
 \tag{E14.1}
\]
This is the local cohomology class under the explicit isomorphism
\[
 V/(F_{E,p}-1)V\ \xrightarrow{\sim}\
 \bigl((V\otimes\mathbb Q_\ell)/V\bigr)^{F_{E,p}},
 \qquad [x]\longmapsto(F_{E,p}-1)^{-1}x+V.
 \tag{E14.2}
\]
The integer $m_p$ is independent of the Frobenius prolongation,
the cocycle representative, the integral Tate basis, and
multiplication of $[h]$ by a unit.

For $q>1$ and $n\ge1$ put
\[
             T_q(n)=\frac{q^{2-n}}{q^2-1}.
\]
The following are natural densities relative to the indicated
set of rational primes; all of them exist:
\[
\begin{array}{c|cc}
 &p\equiv1\pmod4&p\equiv3\pmod4\\\hline
 \ell=3:\ m_p\ge n&T_9(n)&T_3(n)\\
 \ell=5:\ m_p\ge n&2T_5(n)-T_5(n)^2&T_5(n).
\end{array}
 \tag{E14.3}
\]
In particular, among all rational primes, the exact tails are
\[
\begin{split}
 \delta_3(m_p\ge n)&=\frac{9^{2-n}}{160}
                         +\frac{3^{2-n}}{16},\\
 \delta_5(m_p\ge n)&=\frac{5^{2-n}}{16}
                         -\frac{5^{4-2n}}{1152}.
\end{split}
 \tag{E14.4}
\]
The masses at zero and at every positive integer are
\[
\begin{array}{c|cc}
 &\delta(m_p=0)&\delta(m_p=n),\quad n\ge1\\\hline
 \ell=3&121/160&9^{1-n}/20+3^{1-n}/8\\
 \ell=5&817/1152&5^{1-n}/4-5^{2-2n}/48.
\end{array}
 \tag{E14.5}
\]
Each positive exponent has positive density.  For a nonzero scalar
$r\in\mathbb Z_\ell$, with $t=\nu_\ell(r)$, the local exponents
of $r[h]$ are exactly $\max(m_p-t,0)$; their positive tails are
obtained by replacing $n$ with $n+t$ in (E14.3)--(E14.4).
For a compact vector $(r_1[h],r_2[h],r_3[h])$, the analogous
exponent is $\max(m_p-\min_i\nu_\ell(r_i),0)$, with value zero
when every $r_i$ is zero.  These assertions apply to the primitive
class used in the prior E.115 and E.117 and do not require the
prior rank-one cohomology calculation.
\end{statement}
\begin{proof}
First $\det(F_{E,p}-1)=\#E_+(\mathbb F_p)>0$, so every inverse
in (E14.1)--(E14.2) exists.  The inverse to (E14.2) sends a
fixed class $u+V$ to $(F_{E,p}-1)u$ modulo $(F_{E,p}-1)V$.
Both composites are the identity.  Since $h$ is unramified at
$p$, $h(F_p)$ is its unramified cohomology coordinate.

Write the affine group law as
$(x,F)(y,G)=(x+Fy,FG)$.
Conjugation by $(v,B)$ gives
\[
 F'=BFB^{-1},\qquad x'=Bx-(F'-1)v,\qquad
 (F'-1)^{-1}x'=B(F-1)^{-1}x-v.
 \tag{E14.6}
\]
Thus the order of the last vector modulo $V$ is unchanged.
A coboundary adds an integral vector to $(F-1)^{-1}x$, and
a unit scalar preserves its order.  This proves all stated
independences and, in particular, the conjugacy invariance
needed for Chebotarev.

By E.13 the relevant profinite quotient is the full group
$V\rtimes N_\ell$.  Its normalized Haar measure is the product
of additive Haar measure on $V$ and Haar measure on $N_\ell$;
the linear automorphisms preserve the former measure.
The two Cartan cosets have measure one half.  Their arithmetic
quotient is $\operatorname{Gal}(\mathbb Q(i)/\mathbb Q)$, so they
correspond exactly to $p\equiv1$ and $p\equiv3\pmod4$.

Here is the one-coordinate calculation, including its finite
level form.  Let $\mathcal O$ be either $\mathbb Z_\ell$ or
$\mathbb Z_3[i]$, let its residue field have size $q$, and
normalize its valuation by $v(\ell)=1$.  Let $a$ and $x$ be
independent uniform elements of $\mathcal O^\times$ and
$\mathcal O$.  For $j\ge0$ and $k\ge1$,
\[
 \Pr(v(x)=j)=(1-q^{-1})q^{-j},\qquad
 \Pr(v(a-1)\ge k)=\frac{q^{1-k}}{q-1}.
\]
The exponent of $x/(a-1)$ modulo $\mathcal O$ is
$\max(v(a-1)-v(x),0)$ outside a set of measure zero.
For $n\ge1$ its tail is therefore
\[
 \sum_{j\ge0}(1-q^{-1})q^{-j}
          \frac{q^{1-j-n}}{q-1}
   =q^{-n}\sum_{j\ge0}q^{-2j}
   =T_q(n).
 \tag{E14.7}
\]
Modulo $\ell^M$, the exponent of $[x]$ in
$\mathcal O/((a-1),\ell^M)$ has the exact tail
\[
 T_{q,M}(n)=T_q(n)\bigl(1-q^{-2(M-n+1)}\bigr)
       \quad(1\le n\le M),
 \tag{E14.8}
\]
and zero tail for $n>M$.  This is the same sum with
$0\le j\le M-n$.

In the Cartan coset at five, $F=(a,b)$ has two independent
unit coordinates and $x=(x_1,x_2)$ two independent additive
coordinates.  The exponent is the maximum of the two
one-coordinate exponents, giving $2T_5-T_5^2$.
At three the Cartan is the single unit group of the unramified
quadratic ring, so (E14.7) gives $T_9$ directly.

In the other coset write $F=a c$.  Its square is $N(a)$,
its trace is zero and its determinant is $-N(a)$.
Modulo the odd prime $\ell$, $F$ cannot equal the identity:
it anticommutes with the invertible CM operator $I$, and
$2I\ne0$ modulo $\ell$.  At least one entry of $F-1$ is
therefore a unit.  The Smith factors of this matrix over
$\mathbb Z_\ell$ are $1$ and $1-N(a)$, up to units.
For fixed $F$, the image of a uniform $x\in V$ in the
nontrivial Smith coordinate is uniform in $\mathbb Z_\ell$.
Moreover, $N(a)$ is uniform in $\mathbb Z_\ell^\times$.
At five this is the product map on two independent units.
At three, the residue-field norm is onto and the norm on
principal units is onto: the $3$-adic logarithm identifies it
with the trace $\mathbb Z_3[i]\to\mathbb Z_3$, whose value at
$1$ is the unit two.  Haar measure under this surjective
homomorphism is normalized Haar measure.
Consequently the other coset has tail $T_\ell(n)$, and the
same argument modulo $\ell^M$ gives $T_{\ell,M}(n)$.
For the actual arithmetic Frobenius at an inert rational
prime, the determinant identity reads $N(a)=-p$ and
$\det(F-1)=p+1$, retaining the sign in this calculation.

For completeness the passage to prime density uses finite
quotients with a vanishing error, not an interchange of an
uncontrolled infinite union with density.
Let $A_{n,M}$ be the inverse image in $V\rtimes N_\ell$ of
\[
 \ell^{n-1}x\notin(F-1)(V/\ell^M V).
\]
This is a conjugacy-invariant clopen set.  Put
$D_M=\{\det(F-1)\equiv0\pmod{\ell^M}\}$.
For a matrix with $v_\ell(\det(F-1))<M$, both Smith exponents
are less than $M$.  Solvability of the displayed integral
equation is consequently detected modulo $\ell^M$.
Thus, for the actual Frobenius elements,
\[
 A_{n,M}\ \subset\ \{m_p\ge n\}\ \subset\ A_{n,M}\cup D_M.
 \tag{E14.9}
\]
The measures of $D_M$ tend to zero.  In the split Cartan,
the limiting set is $a=1$ or $b=1$; in the nonsplit Cartan
it is $a=1$; and in the other coset it is $N(a)=1$.
Each has measure zero by the product unit measures or the
surjective norm map just proved.  Finite-level Chebotarev
in its natural-density form applies to both clopen sets in
(E14.9) \cite[\S2.1, Theorem 1, and equations (8)--(9)]{E_density_Serre}.
Taking lower and upper natural densities and then $M\to\infty$
proves existence and the asserted Haar values.  The same
sandwich intersected with a Cartan coset proves the relative
densities in (E14.3).

Averaging the two cosets yields (E14.4); subtracting successive
tails gives the positive masses in (E14.5), and subtracting the
first tail from one gives its masses at zero.
Multiplication by $r$ on (E14.2) gives the stated exponent
shift.  For a vector of three scalar multiples, the order in
the direct sum is the maximum of the three scalar orders,
which gives the minimum valuation formula.

The independent finite certificate
\texttt{code/B/cm\_density\_checks.py} enumerates every matrix
in both normalizer cosets and every translation vector modulo
$3$, $3^2$, $3^3$, $5$, and $5^2$.  It forms each actual image
$(F-1)(V/\ell^M V)$ and counts the order of every cokernel
class.  Every count agrees exactly with (E14.8), its
two-coordinate maximum, and the other-coset formula.  Its
output is retained in
\texttt{code/B/cm\_density\_checks.out.json}.
\end{proof}
